\section{Conclusion}
\label{sec:conclusion}

% We presented \textbf{GIFSplat}, an \emph{iterative feed-forward} framework for 3D Gaussian Splatting prediction that bridges the gap between fast one-shot predictors and slow per-scene optimization.
% Instead of test-time backpropagation, GIFSplat performs a short sequence of \emph{forward-only residual updates} that refine gaussian parameters using evidence from render and observation differences.
% A \emph{generative prior fusion} module converts diffusion-enhanced renderings into cues, providing prior-aware corrections without view explosion or heavy optimization loops. Extensive experiments on DL3DV, DTU and RealEstate10K show consistent gains over strong feed-forward baselines, with particularly notable improvements in perceptual quality.
% Against optimization-based generative pipelines, GIFSplat attains superior reconstruction quality while avoiding costly test-time optimization in sparse view setting. Ablations confirm the complementarity of iterative residual updates and prior fusion.


We presented GIFSplat, an iterative feed-forward 3D Gaussian Splatting framework that refines an initial prediction through multi-step forward-only residual updates guided by both observations and diffusion-enhanced cues. Our IFSplat baseline already brings scene-specific refinement without test-time backpropagation, and the full GIFSplat variant further injects a frozen generative prior via lightweight Gaussian-level discrepancy cues. Experiments on DL3DV, RealEstate10K, and DTU show consistent gains over recent feed-forward 3D reconstruction methods, especially in cross-domain settings.


\noindent\textbf{Limitations.}
% The transition from pixel-form interactions to point-based interactions in our work still leaves room for even more performant designs: an ideal model would first aggregate 2D features into 3D features, perform all interactions in 3D attention space, and then directly predict 3D Gaussians, rather than aggregating pixel-aligned Gaussians and only afterwards refining them with 3D attention. While such a fully point-based formulation could further improve performance, it is orthogonal to the main focus of this work. Our primary knowledge contribution is to show that an iterative residual scheme can progressively enhance 3D reconstruction quality in a purely feed-forward manner, while enabling a generative prior to be injected into this process in a simple yet effective way. We hope GIFSplat provides a practical design point for scalable, high-quality, and robust feed-forward 3D reconstruction.
Our refinement head still focus on static scenes and a small, fixed set of input modalities. Extending GIFSplat to dynamic content or to flexible geometric priors such as known depth maps or normal maps is an interesting direction for future work.

%\paragraph{References}
%\cite{mildenhall2021nerf,muller2022instant,kerbl20233d,wang2024dust3r,leroy2024grounding,yang2025fast3r,wang2025vggt,zhang2025flare,ye2024no,jiang2025anysplat,charatan2024pixelsplat,chen2024mvsplat,xu2025depthsplat,smart2024splatt3r,keetha2025mapanything,wang2025pi,poole2022dreamfusion,yu2024viewcrafter,wang2024splatflow,wu2025genfusion,wu2025difix3d+,fischer2025flowr,guo2023ree3d,liu20243dgs,kang2025ilrm,anftl2021vision,zhang2018unreasonable,wang2004image,zhou2018stereo,ling2024dl3dv}


% 3D Reconstruction: 3DGS, Mip-Splatting, InstantSplat

% Generative 3R: 
% opensource: ViewCrafter,  Genfusion, Difix3D+
% closesource: FlowR, Wonderland, GS-Enhancer, GSfix3D

% GFMs: Dust3R, Mast3R, Splatt3R, Fast3R, VGGT, MapAnything

% Feed-forward: PixelSplat, MVSplat, DepthSplat, NopoSplat, FLARE, AnySplat


% [1] Mildenhall et al., “NeRF: Representing Scenes as Neural Radiance Fields for View Synthesis,” ECCV 2020.

% [2] Kerbl et al., “3D Gaussian Splatting for Real-Time Radiance Field Rendering,” SIGGRAPH 2023.

% [3] Müller et al., “Instant Neural Graphics Primitives with a Multiresolution Hash Encoding,” SIGGRAPH 2022.

% [4] Wang et al., “DUSt3R: Geometric 3D Vision Made Easy,” CVPR 2024.

% [5] Li et al., “Grounding Image Matching in 3D with MASt3R,” arXiv 2024.

% [6] Wang et al., “VGGT: Visual Geometry Grounded Transformer,” CVPR 2025.

% [7] Revaud et al., “Fast3R: Towards 3D Reconstruction of 1000+ Images in One Forward Pass,” CVPR 2025.

% [8] Yu et al., “FLARE: Feed-forward Geometry, Appearance and Camera Estimation from Uncalibrated Sparse Views,” CVPR 2025.

% [9] Chen et al., “NoPoSplat: Learning to Splat without Poses,” ICCV 2023.

% [10] Jiang et al., “AnySplat: Feed-forward 3D Gaussian Splatting from Unconstrained Views,” SIGGRAPH Asia 2025.

% [11] Poole et al., “DreamFusion: Text-to-3D using 2D Diffusion,” arXiv 2022.

% [12] Chen et al., “ViewCrafter: Taming Video Diffusion Models for High-Fidelity Novel View Synthesis,” arXiv 2024.

% [13] Go et al., “SplatFlow: Multi-View Rectified Flow Model for 3D Gaussian Splatting Synthesis,” CVPR 2025.

% [14] Wu et al., “GenFusion: Closing the Loop between Reconstruction and Generation via Videos,” CVPR 2025.

% [15] Liu et al., “3DGS-Enhancer: Enhancing Unbounded 3D Gaussian Splatting with View-consistent 2D Diffusion Priors,” NeurIPS 2024.

% [16] Wu et al., “Difix3D+: Improving 3D Reconstructions with Single-Step Diffusion Models,” CVPR 2025.

% [17] Fischer et al., “FlowR: Flowing from Sparse to Dense 3D Reconstructions,” ICCV 2025.